\par\Needspace{6\baselineskip}\medskip
\begin{center}\hypertarget{native-172-article125}{}\pdfbookmark[2]{125.}{article-125}125.\end{center}
Levi mutatione applicata e methodo praecedente \emph{secundam} eliciemus. A
distantiis in duabus observationibus profecti, perinde ut in illa, cuncta elementa
determinabimus; ex his vero non locum geocentricum pro observatione
tertia computabimus, sed tantummodo usque ad locum heliocentricum in orbita
progrediemur; ex altera parte eundem locum heliocentricum per problema
in artt. 74, 75 tractatum e loco geocentrico observato atque situ plani orbitae
derivabimus; hae duae determinationes inter se differentes (nisi forte valores
veri ipsarum $x,y$ suppositi fuerint) ipsas $X,Y$ nobis suppeditabunt, accepta
pro $X$ differentia inter duos valores longitudinis in orbita, atque pro $Y$ \SourceLeaf{173}differentia
inter duos valores radii vectoris, aut potius logarithmi eius. Haecce
methodus iisdem monitionibus obnoxia est, quas in art. praec. attigimus: adiungere
oportet aliam, scilicet, quod locus heliocentricus in orbita e geocentrico
deduci nequit, quoties locus terrae in alterutrum nodorum orbitae incidit;
tunc itaque hanc methodum applicare non licet. Sed in eo quoque
casu, ubi locus terrae ab alterutro nodorum perparum distat, hac methodo abstinere
conveniet, quoniam suppositio, variationibus parvis ipsarum $x,y$ respondere
variationes proportionales ipsarum $X,Y$, nimis erronea evaderet, per
rationem ei quam in art. praec. attigimus similem. Sed hic quoque remedium
e permutatione loci medii cum aliquo extremorum, cui locus terrae a nodis
magis remotus respondeat, petere licebit, nisi forte in omnibus tribus observationibus
terra in nodorum viciniis versata fuerit.
